\begin{frame}
	\frametitle{A camada BitLinear}
	\only<1>{
		\framesubtitle{Da camada Linear densa à ternária}
		\par Uma camada Linear comum calcula $Y = W X$, com $W$ em FP16/FP32: cada saída soma \textbf{multiplicações reais} entre pesos e ativações contínuos.\newline
		
		\par O BitNet b1.58 \cite{ma2024era} substitui $W$ por uma versão ternária $\widetilde{W} \in \{-1,0,1\}$, obtida pela \textbf{quantização \textit{absmax}}:
		
		\begin{equation}
			\gamma = \text{mean}(|W|), \qquad
			\widetilde{W} = \text{round}\big(\text{clip}(W/\gamma,\,-1,\,1)\big)
		\end{equation}
		\par $\gamma$ é um único número por tensor --- guardado à parte para reescalar a saída no fim, não para cada peso individualmente.
	}
	\only<2>{
		\framesubtitle{Exemplo numérico}
		\par Tome três pesos de uma linha de $W$: $w = [0{,}6,\ -0{,}3,\ 0{,}05]$.
		\begin{equation*}
			\gamma = \frac{0{,}6+0{,}3+0{,}05}{3} = \frac{0{,}95}{3} = 0{,}3167
		\end{equation*}
		% quebrado em duas linhas: os tres estagios numa linha so estouravam a
		% largura do slide (Overfull \hbox de 14,5pt)
		\begin{equation*}
			\begin{aligned}
				w/\gamma &= [1{,}895,\ -0{,}947,\ 0{,}158]
				\xrightarrow{\text{clip}(-1,1)} [1{,}0,\ -0{,}947,\ 0{,}158] \\
				&\xrightarrow{\text{round}} \boxed{[1,\ -1,\ 0]}
			\end{aligned}
		\end{equation*}
		\par O peso original $0{,}05$ (o menor em módulo) \textbf{vira zero} --- é descartado, não aproximado. Os outros dois viram $\pm 1$.
		\par As ativações $X$ recebem seu próprio tratamento, quantização \textbf{absmax} para INT8: $\widetilde{X} = \text{round}(X \cdot 127/\max(|X|))$, com \textit{clip} em $[-128,127]$.
		% [SOFTWARE] BitNet -> Quantização | slider w ~ 0 para reproduzir a
		% zona morta, depois 0,6 / -0,3 para os dois outros pesos do exemplo
		\abrirNoSoftware{BitNet \textrightarrow{} Quantização}{bitnet.quant}
	}
	\only<3>{
		\framesubtitle{O que essa troca faz}
		\begin{itemize}
			\item Com $\widetilde{W} \in \{-1,0,1\}$, a multiplicação $\widetilde{W}\cdot\widetilde{X}$ vira: \newline 
			\textbf{some} $\widetilde{X}$ (peso $+1$), \newline
			\textbf{subtraia} $\widetilde{X}$ (peso $-1$), ou \newline
			\textbf{ignore} (peso $0$).\newline
			
			\item Nenhuma multiplicação real acontece no lado dos pesos --- só somas/subtrações de inteiros, mais um único rescalonamento por $\gamma$ ao final.
		\end{itemize}
	}
\end{frame}
